% Lecture notes.
% \lecture{date}{title} starts a new lecture. Definitions, theorems and examples
% are numbered within the lecture and drawn as coloured boxes (tcolorbox).
\documentclass[11pt,a4paper]{article}

\usepackage[a4paper,margin=2.5cm,marginparwidth=2cm]{geometry}
\usepackage[english]{babel}
\usepackage{microtype}
\usepackage{amsmath,amssymb,amsthm}
\usepackage[dvipsnames]{xcolor}
\usepackage[most]{tcolorbox}
\usepackage{enumitem}
\usepackage[hidelinks]{hyperref}

% ---- Course details: change these ------------------------------------------
\title{Course Name\\[0.3em]\large Lecture Notes}
\author{Lecturer Name}
\date{Term 2026}

% ---- Lectures -----------------------------------------------------------------
\newcounter{lecture}
\newcommand{\lecture}[2]{%
  \refstepcounter{lecture}%
  \section*{Lecture \thelecture: #2}%
  \addcontentsline{toc}{section}{Lecture \thelecture: #2}%
  \marginpar{\raggedright\footnotesize\color{gray}#1}}
\counterwithin*{equation}{lecture}

% ---- Boxes ----------------------------------------------------------------------
% Change the colours or add your own with \newtcbtheorem.
\tcbset{notebox/.style={enhanced,breakable,boxrule=0.6pt,arc=2pt,
  fonttitle=\bfseries,separator sign={.},description font=\mdseries\itshape}}
\newtcbtheorem[number within=lecture]{definition}{Definition}%
  {notebox,colback=NavyBlue!5,colframe=NavyBlue!70!black}{def}
\newtcbtheorem[use counter from=definition]{theorem}{Theorem}%
  {notebox,colback=Plum!5,colframe=Plum!70!black}{thm}
\newtcbtheorem[use counter from=definition]{example}{Example}%
  {notebox,colback=ForestGreen!5,colframe=ForestGreen!70!black}{ex}
\newcommand{\remark}[1]{\marginpar{\raggedright\footnotesize\color{BrickRed}#1}}

\begin{document}
\maketitle
\tableofcontents

\lecture{6 Jan 2026}{Sets and functions}

\begin{definition}{Function}{function}
  A \emph{function} $f\colon A \to B$ assigns to every $a \in A$ exactly one
  element $f(a) \in B$.
\end{definition}

A function is \emph{injective} if different inputs give different outputs.%
\remark{Also called one-to-one.}

\begin{example}{}{square}
  $f\colon \mathbb{R} \to \mathbb{R}$, $f(x) = x^2$ is not injective, since
  $f(-1) = f(1)$.
\end{example}

\lecture{8 Jan 2026}{Limits}

\begin{definition}{Limit}{limit}
  We write $\lim_{x \to a} f(x) = L$ if for every $\varepsilon > 0$ there is a
  $\delta > 0$ such that $0 < |x - a| < \delta$ implies $|f(x) - L| < \varepsilon$.
\end{definition}

\begin{theorem}{Sum rule}{sum}
  If $\lim_{x\to a} f(x) = L$ and $\lim_{x\to a} g(x) = M$, then
  \begin{equation}
    \lim_{x\to a} \bigl(f(x) + g(x)\bigr) = L + M.
  \end{equation}
\end{theorem}

\begin{proof}
  Given $\varepsilon > 0$, choose $\delta$ so that both $|f(x) - L|$ and
  $|g(x) - M|$ are below $\varepsilon/2$; the triangle inequality does the rest.
\end{proof}

Theorem~\ref{thm:sum} relies on Definition~\ref{def:limit}.

\subsection*{Exercises}
\begin{enumerate}[label=\arabic*.]
  \item Show that $f(x) = 2x + 1$ is injective.
  \item Use Definition~\ref{def:limit} to prove that $\lim_{x\to 2} 3x = 6$.
\end{enumerate}

\end{document}
